% Generated draft: requires mathematical and source review.
\documentclass[11pt]{article}
\usepackage[margin=25mm]{geometry}
\usepackage{fontspec}
\setmainfont{Latin Modern Roman}
\usepackage{amsmath,amssymb,mathtools,unicode-math}
\setmathfont{Latin Modern Math}
\usepackage{enumitem,needspace,xurl,hyperref}
\setlength{\parindent}{0pt}
\setlength{\parskip}{5pt}
\setlength{\emergencystretch}{3em}
\newcommand{\N}{\mathbb{N}}
\newcommand{\Z}{\mathbb{Z}}
\newcommand{\Q}{\mathbb{Q}}
\newcommand{\R}{\mathbb{R}}
\newcommand{\C}{\mathbb{C}}
\newcommand{\F}{\mathbb{F}}
\newcommand{\A}{\mathbb{A}}
\newcommand{\PP}{\mathbb{P}}
\newcommand{\E}{\mathbb{E}}
\newcommand{\eps}{\varepsilon}
\newcommand{\dd}{\,\mathrm{d}}
\begin{document}
\section*{UnsolvedMath ID: 20002902}
\textbf{Source problem code:} catalogue entry
\subsection{Definitions and mathematical statement}
The workshop asks to elucidate the representation-theoretic significance of $k$-Schur functions, allowing either symmetric-group or Lie-algebra representation theory. Fix $k\geq1$, a $k$-bounded partition $\lambda$ with $\lambda_1\leq k$, and $m=|\lambda|$. Write $s_\lambda^{(k)}(X;u)$ for the modern graded $k$-Schur function in the normalization of the constructions reported below. The original record does not define its candidate or parameter convention, so identification with every historical $k$-Schur definition is not assumed.

The supplied aid reports Lam's published identification
\[
\C[h_1,\ldots,h_k]\simeq H_*(\operatorname{Gr}_{\operatorname{SL}_{k+1}};\C),\qquad
s_\lambda^{(k)}(X;1)\longmapsto\xi_{w_\lambda},
\]
where $h_j$ are complete symmetric functions, the target is affine-Grassmannian homology with Pontryagin product, and $\xi_{w_\lambda}$ is the Schubert class indexed by the affine Grassmannian element corresponding to $\lambda$. Symmetric degree $m$ corresponds to homological degree $2m$.

It further reports that Kato's current preprint constructs a finite-dimensional graded module $M_\lambda^{(k)}$ over
\[
\C[S_m]\ltimes\C[x_1,\ldots,x_m],
\]
with $S_m$ permuting the polynomial variables, together with a finite-dimensional current-algebra partner via affine Schur--Weyl duality. Its graded Frobenius characteristic realizes the full modern function:
\[
\operatorname{Frob}_u(M_\lambda^{(k)})=s_\lambda^{(k)}(X;u).
\]
For a graded $S_m$-module $M=\bigoplus_dM_d$ and Specht module $L_\mu$, the convention here is
\[
\operatorname{Frob}_u(M)=\sum_{\mu\vdash m}\sum_d
u^d\dim_\C\operatorname{Hom}_{S_m}(L_\mu,M_d)\,s_\mu(X),
\]
The actual polynomial grading and any duality or shift conventions must be matched to the preprint.

The aid's bridge asserts that forgetting grading gives exactly Lam's symmetric-function representative:
\[
[s_\mu]s_\lambda^{(k)}(X;1)
=\sum_d\dim_\C\operatorname{Hom}_{S_m}(L_\mu,(M_\lambda^{(k)})_d).
\]
Thus its Schur coefficients are total Specht multiplicities. Ordinary homology records the specialization $u=1$ and cannot recover these graded multiplicities without an enhancement.

For $k\geq m$, the aid additionally asserts that the grading-forgotten module is the simple Specht module $L_\lambda$ with every $x_i$ acting by zero, and that its current-algebra partner is the evaluation module $V_\lambda^*$ at $z=0$. An evaluation module means that $a\otimes f(z)$ acts as $f(0)a$; the Lie algebra and rank range of this construction are not specified in the source record. The summary mentions an unstable example at $k=2$, $m=3$, but gives neither its module nor its graded character, so no missing decomposition is invented. These claims illustrate the original representation-theoretic request subject to the cited constructions' omitted hypotheses.
\subsection{Short English statement}
The supplied aid interprets modern graded $k$-Schur functions as characters of symmetric-group/current-algebra modules and their parameter-one specialization as affine-Grassmannian Schubert classes. Its bridge identifies Schur coefficients with total Specht multiplicities and separates graded information from ordinary homology.
\subsection{Sources}
\begin{itemize}
\item Local UnsolvedMath registry, permanent ID $20002902$, source code \texttt{catalogue entry}; canonical \texttt{aim-special-functions-notes.json}, \texttt{notes[6]}.
\item AIM workshop \emph{Generalized Kostka polynomials}, Problem (7): \url{https://aimath.org/WWN/kostka/prob2.pdf}. The source explicitly allows symmetric-group or Lie-algebra interpretations; no original remarks or literature field are supplied.
\item The research summary names Lam's published affine-Grassmannian homology theorem and Kato's current preprint on graded $k$-Schur/current-algebra modules, without titles, dates, identifiers or full authoritative citations. The difference in publication status is retained.
\item Research attempt 1: \texttt{partial\_result}, status \texttt{solved}, classification \texttt{SOLVED-IN-LITERATURE}, rigor \texttt{proved}, novelty confidence low. Its bridge, stability and example claims are machine-generated aid, not peer reviewed, and require independent verification before citation.
\end{itemize}
\end{document}
